\section*{Hoofdstuk \ref{ch:b2:anthropogenic-climate-change} --- Antropogene klimaatverandering}

\begin{solution}{exo:b2:anthropogenic-climate-change:1}
(a) $3.75 - 0.60 + 0.22 + 0.45 + 0.11 = \qty{3.93}{W.m^{-2}}$, waarbij de aanpassingen samen
\qty{0.18}{W.m^{-2}} bijdragen. (b) $3.93/3.71 = 1.059$: \qty{5.9}{\percent} meer. (c) De aanpassing van de
troposferische temperatuur, $-0.60$: een warmere troposfeer zendt meer infrarood uit naar de ruimte, wat
een deel van de extra energie wegneemt voordat het oppervlak heeft gereageerd.
\end{solution}

\begin{solution}{exo:b2:anthropogenic-climate-change:2}
(a) $4\pi(\num{6.371e6})^2 = \qty{5.10e14}{m^2}$. (b) Over 47 jaar: $\num{434.9e21}/(47 \times
\num{3.156e7} \times \num{5.10e14}) = \qty{0.575}{W.m^{-2}}$; over 12 jaar: $\num{152.4e21}/(12
\times \num{3.156e7} \times \num{5.10e14}) = \qty{0.789}{W.m^{-2}}$, de 0.57 en 0.79 van de
beoordeling. (c) $0.91 \times 434.9 = \qty{396}{ZJ}$.
\end{solution}

\begin{solution}{exo:b2:anthropogenic-climate-change:3}
(a) $-3.22 + 1.30 + 0.35 + 0.42 - 0.01 = \num{-1.16}$, de netto waarde. (b) $3.93/1.16 =
\qty{3.39}{\celsius}$; $3.93/3.22 = \qty{1.22}{\celsius}$. (c) $3.39/1.22 = 2.78$: de terugkoppelingen
verdrievoudigen de Planck-opwarming bijna.
\end{solution}

\begin{solution}{exo:b2:anthropogenic-climate-change:4}
(a) Eén petagram koolstof is $44.009/12.011 = 3.664$~Gt koolstofdioxide: $0.45 \times 3.664 =
\qty{1.65}{\celsius}$ per \qty{1000}{PgC}, en het bereik wordt 0.989 tot 2.31. (b) $2390 \times
12.011/44.009 = \qty{652}{PgC}$. (c) $0.45 \times 2.39 = \qty{1.08}{\celsius}$.
\end{solution}

\begin{solution}{exo:b2:anthropogenic-climate-change:5}
(a) $g = 159/398 = 0.400$; $T = (398/\num{5.670e-8})^{1/4} = \qty{289}{K}$. (b) $-0.6 \times 4
\times \num{5.670e-8} \times 288^3 = \qty{-3.25}{W.m^{-2}.K^{-1}}$; bij \qty{298}{K} $-3.60$.
(c) Minder gevoelig: de respons groeit als $T^3$, dus een warmere planeet raakt elke extra graad
gemakkelijker kwijt, bij gelijke andere terugkoppelingen.
\end{solution}

\begin{solution}{exo:b2:anthropogenic-climate-change:6}
(a) $1.85 - 0.50 = \qty{1.35}{W.m^{-2}.\celsius^{-1}}$. (b) $\sqrt{0.32^2 + 0.20^2} = 0.38$.
(c) De som is veel beter bekend dan elk van de termen: hun fouten zijn sterk anticorreleerd. Waar
de bovenste troposfeer meer opwarmt, bevat ze ook meer damp, dus een sterkere negatieve \omterm{def:b2:anthropogenic-climate-change:feedbacks}{terugkoppeling van de
temperatuurgradiënt} gaat samen met een sterkere positieve \omterm{def:b2:anthropogenic-climate-change:feedbacks}{waterdampterugkoppeling}.
\end{solution}

\begin{solution}{exo:b2:anthropogenic-climate-change:7}
(a) $\lambda = (0.59 - 2.20)/1.03 = \qty{-1.56}{W.m^{-2}.\celsius^{-1}}$; $3.93/1.56 =
\qty{2.52}{\celsius}$. (b) $3.93/(1.56 - 0.5) = 3.93/1.06 = \qty{3.71}{\celsius}$. (c) $3.93/1.46 =
\qty{2.69}{\celsius}$. De beoordeling vindt met volledige verdelingen 3.5 en \qty{2.7}{\celsius}.
\end{solution}

\begin{solution}{exo:b2:anthropogenic-climate-change:8}
(a) $0.617/3682 = \qty{1.68e-4}{\celsius^{-1}}$. (b) $0.113 \times 0.396 = \qty{0.0447}{m}$, of
\qty{44.7}{mm}, dicht bij de \qty{47.5}{mm} van het budget. (c) $\num{2e-4} \times 0.2 \times 700 =
\qty{0.028}{m}$, \qty{28}{mm}.
\end{solution}

\begin{solution}{exo:b2:anthropogenic-climate-change:9}
(a) $\sum xy = 0.311$, $\sum x^2 = 0.300$: $\beta = 1.04$. (b) Het interval van \qty{95}{\percent} is
$1.04 \pm 1.96 \times 0.08$, 0.88 tot 1.20: het sluit 0 uit, dus de respons is gedetecteerd, en het bevat
1, dus de waargenomen verandering is in overeenstemming met de gesimuleerde. (c) Het interval 0.34 tot 0.66 sluit
0 nog steeds uit (gedetecteerd), maar sluit 1 uit: het model zou de respons ongeveer een factor
twee overschatten.
\end{solution}

\begin{solution}{exo:b2:anthropogenic-climate-change:10}
(a) $\sum g^2 = 1.20$, $\sum a^2 = 0.18$, $\sum ga = -0.46$, $\sum gy = 0.74$, $\sum ay = -0.28$.
De normaalvergelijkingen zijn $1.20\,\beta_g - 0.46\,\beta_a = 0.74$ en $-0.46\,\beta_g + 0.18\,\beta_a
= -0.28$; hun determinant is $1.20 \times 0.18 - 0.46^2 = 0.0044$, en $\beta_g = (0.74 \times
0.18 - 0.46 \times 0.28)/0.0044 = 1.00$, $\beta_a = (-1.20 \times 0.28 + 0.46 \times 0.74)/0.0044 =
1.00$. (b) Broeikasgassen $+\qty{0.80}{\celsius}$, aerosolen $-\qty{0.30}{\celsius}$. (c) De
determinant is klein omdat $g$ en $a$ bijna evenredig zijn (hun correlatie is $-0.99$): een
beetje ruis in $y$ kan worden opgevangen door beide schaalfactoren samen te verhogen, zodat elk slecht
bepaald is, al is hun gecombineerde effect goed bekend.
\end{solution}

\begin{solution}{exo:b2:anthropogenic-climate-change:11}
(a) $(0.93 - 0.1)/0.45 \times 1000 = \qty{1844}{Gt}$; min 210 en $26 \times 0.93 = 24$:
\qty{1610}{Gt}. (b) $0.93 - x = 0.45 \times (1350 + 210 + 24)/1000 = 0.713$, dus $x =
\qty{0.22}{\celsius}$: ongeveer de bovenkant van het beoordeelde bereik 0.1 tot \qty{0.2}{\celsius}. (c) Het
budget hangt van de opwarming door andere factoren dan CO$_2$, $x$, alleen af via $-x/\text{TCRE}$: elke
\qty{0.1}{\celsius} kost $0.1/0.45 \times 1000 = \qty{222}{Gt}$, wat de grens ook is.
\end{solution}

\begin{solution}{exo:b2:anthropogenic-climate-change:12}
(a) $3.93/1.81 = \qty{2.17}{\celsius}$, $3.93/0.51 = \qty{7.71}{\celsius}$, $3.93/1.16 =
\qty{3.39}{\celsius}$. (b) De lage kant is $3.39 - 2.17 = \qty{1.22}{\celsius}$ lang, de hoge kant
$7.71 - 3.39 = \qty{4.32}{\celsius}$: de functie $1/|\lambda|$ is convex, dus gelijke stappen van
$\lambda$ naar nul verhogen de ECS steeds meer. (c) Een zwak negatieve $\lambda$ geeft een zeer
grote ECS, en kleine veranderingen van $\lambda$ dicht bij nul zijn moeilijk uit te sluiten; een lage ECS vraagt een sterk
negatieve $\lambda$, die door de \omterm{def:b2:anthropogenic-climate-change:feedbacks}{Planck-respons} wordt begrensd en door de waargenomen opwarming wordt tegengesproken.
\end{solution}

\begin{solution}{pb:b2:anthropogenic-climate-change:1}
\textbf{1.} $2.16 + 0.54 + 0.21 + 0.41 = \qty{3.32}{W.m^{-2}}$.
\textbf{2.} $-0.22 - 0.84 = \qty{-1.06}{W.m^{-2}}$.
\textbf{3.} $3.32 + 0.47 + 0.05 - 1.06 - 0.20 + 0.08 + 0.06 = \qty{2.72}{W.m^{-2}}$, het beoordeelde
totaal.
\textbf{4.} $3.93\log_2(410/278) = \qty{2.20}{W.m^{-2}}$, tot op \qty{2}{\percent} gelijk aan 2.16.
\textbf{5.} $3.93\log_2(427.35/278) = \qty{2.44}{W.m^{-2}}$.
\textbf{6.} De term aerosol--wolken, $-1.45$ tot $-0.25$, \qty{1.20}{W.m^{-2}} breed, tegenover 1.52 voor
het totaal: hij draagt het grootste deel van de onzekerheid van het totaal omdat hij groot en slecht
bepaald is.
\textbf{7.} $0.79 \times \num{5.10e14} \times 12 \times \num{3.156e7} = \qty{1.53e23}{J}$,
\qty{153}{ZJ}.
\textbf{8.} $0.91 \times 153 = \qty{139}{ZJ}$.
\textbf{9.} $\lambda = (0.59 - 2.20)/1.03 = \qty{-1.56}{W.m^{-2}.\celsius^{-1}}$.
\textbf{10.} $3.93/1.56 = \qty{2.52}{\celsius}$.
\textbf{11.} $3.93 \times 1.03/2.20 = \qty{1.84}{\celsius}$.
\textbf{12.} $3.93/1.16 = \qty{3.39}{\celsius}$, hoger met \qty{0.87}{\celsius}.
\textbf{13.} De opwarming tot nu toe was het sterkst waar de terugkoppelingen het meest negatief zijn; naarmate het
patroon verandert, wordt de terugkoppelingsparameter minder negatief, zodat de gevoeligheid op lange termijn
de effectieve overtreft.
\textbf{14.} $0.45 \times 44.009/12.011 = \qty{1.65}{\celsius}$ per \qty{1000}{PgC}.
\textbf{15.} 0.989 tot \qty{2.31}{\celsius} per \qty{1000}{PgC}.
\textbf{16.} $2390 \times 12.011/44.009 = \qty{652}{PgC}$.
\textbf{17.} $0.45 \times 2.39 = \qty{1.08}{\celsius}$.
\textbf{18.} Binnen de afronding verklaart de koolstofdioxide alleen de hele menselijke opwarming: de
opwarming door de andere broeikasgassen en de afkoeling door de aerosolen heffen elkaar bijna op.
\textbf{19.} $1.5 - 1.07 = \qty{0.43}{\celsius}$.
\textbf{20.} $(0.43 - 0.1)/0.45 \times 1000 = \qty{733}{Gt}$ koolstofdioxide.
\textbf{21.} $733 - 210 = \qty{523}{Gt}$.
\textbf{22.} $26 \times 0.43 = \qty{11.2}{Gt}$; $523 - 11.2 = \qty{512}{Gt}$.
\textbf{23.} Het daalt met $0.1/0.45 \times 1000 = \qty{222}{Gt}$, tot \qty{290}{Gt}: de
\qty{220}{Gt} waarmee volgens de beoordeling de keuzes over de andere factoren het budget verschuiven.
\textbf{24.} Berekend uit de TCRE is het \omterm{def:b2:anthropogenic-climate-change:budget}{resterende koolstofbudget} voor \qty{1.5}{\celsius} vanaf het
begin van 2020 \textbf{\qty[text-series-to-math=true]{512}{Gt}} koolstofdioxide, tot op
\qty{3}{\percent} gelijk aan de beoordeelde \qty{500}{Gt} bij een waarschijnlijkheid van \qty{50}{\percent}.
\end{solution}
